\documentclass[12pt,A4paper]{article}

\usepackage{times}
\usepackage{amssymb}
\usepackage{amsmath}
\usepackage[margin=1.1in]{geometry}
\usepackage{graph}
\usepackage{graphicx}
\usepackage{hyperref}
%\usepackage{subfig}
\usepackage{pdflscape}
\usepackage{adjustbox}
\usepackage{setspace}
\usepackage{cellspace}


\usepackage{enumitem}
\usepackage[utf8]{inputenc}


 \usepackage{graphicx}
\usepackage{array}
\usepackage{geometry}
\usepackage{mathrsfs}

\usepackage{color}
\usepackage{amsthm}
%\usepackage{enumerate}
\newtheorem{Theorem}{Theorem}
\newtheorem{Definition}{Definition}
\newtheorem{Lemma}{Lemma}
\newtheorem{Corollary}{Corollary}
\newtheorem{Remark}{Remark}
\newtheorem{Observation}{Observation}
\newtheorem{Claim}{Claim}
\newtheorem{Fact}{Fact}
\def\bibsection{\section*{References}}
\newcommand{\com}[1]{}
\newcommand{\boxi}{\mathrm{box}}
\newcommand{\tw}{\mathrm{tw}}
\newcommand{\pw}{\mathrm{pw}}

\newenvironment{claim}{\vspace{0.1in}\par\noindent\textit{Claim. }}
{\vspace{0.1in}\par}

\onehalfspacing
\setlength\cellspacetoplimit{3pt}
\setlength\cellspacebottomlimit{3pt}
\renewcommand{\figurename}{Fig.}
\renewcommand{\contentsname}{\fontsize{20pt}{1}\selectfont Contents}

\title{Give your thesis title}
\date{}
\author{Name in caps}

\setcounter{secnumdepth}{3}
\setcounter{tocdepth}{3}


%\newpage
\pagenumbering{roman} 
\setcounter{page}{-1}

\begin{document}

\begin{center}
\vspace*{1cm}
\textbf{\fontsize{11pt}{1}\selectfont \underline{SYNOPSIS OF THESIS}}
\end{center}

\vspace{0.35cm}
\begin{center}
\fontsize{16pt}{18}\selectfont \textbf{$<$TITLE OF THE THESIS$>$}
\end{center}



\vspace{0.25cm}

 
\begin{center}
\textit{Submitted in partial fulfillment of the\\
requirements for the award of the degree of}
\end{center}
\vspace{-0.5cm}
\begin{center}
\textbf{DOCTOR OF PHILOSOPHY}
\end{center}

\vspace{-0.35cm}
\fontsize{11pt}{14}\selectfont{
\begin{center}
\textit{by}
\end{center}}

\vspace{-0.35cm}
\fontsize{11pt}{14}\selectfont{
\begin{center}
\textbf{$<$NAME OF THE CANDIDATE$>$\\}
\textbf{ROLL NO: $<$PXXXXXXXX$>$}
\end{center}}




\vspace*{0.75cm}
\begin{figure}[htb]
\centering
\includegraphics[trim = 1mm 5mm 1mm 3mm, width=0.3\textwidth]{nitc-logo.png}
\end{figure}

\vspace*{0.5cm}
\begin{center}
\begin{spacing}{1.25}
 {\textsc{$<$Name of The Department$>$}}\\
\textbf{\textsc{NATIONAL INSTITUTE OF TECHNOLOGY CALICUT} }\\
NIT CAMPUS P.O., KOZHIKODE, 
KERALA  - 673601, INDIA \\
\vspace{15pt}
% \bigskip
$<$ Month $ \& $ Year of Submission $>$ % eg: February 2023
\end{spacing}
\end{center}

\pagenumbering{roman} 
\setcounter{page}{-1}



\doublespacing
\tableofcontents
\onehalfspacing
\newpage

\setcounter{page}{0}
\pagenumbering{arabic}



%------------------------------------------------------------------------------------------------------------


\section{Introduction}

The P versus NP problem is to determine whether every language accepted by some non deterministic algorithm in polynomial time is also accepted by some (deterministic) algorithm in polynomial time. To define the problem precisely it is necessary to give a formal model of a computer. The standard computer model in computability theory is the Turing machine, introduced by Alan Turing in 1936~\cite{turing1936computable}...


\subsection{Brief Literature Review}
This section discusses a brief history of the problem we have addressed in this thesis. The theory of NP-completeness has its roots in computability theory, which
originated in the work of Turing, Church, Godel, and others in the 1930s~\cite{church1936note}. The importance of the P vs NP question stems from the successful theories of NP-completeness~\cite{cook1971complexity}, and complexity-based cryptography, as well as the potentially stunning practical consequences of a constructive proof of P=NP...

\subsection{Motivation/Problem Statement/Objectives}
Many hundreds of NP-complete problems have been identified, including Subset Sum (given a set of positive integers presented in decimal notation, and a target T, is there a subset summing to T?), many graph problems (given a graph G, does G have a Hamiltonian cycle? Does G have a clique consisting of half of the vertices? Can the vertices of G be colored with three colors with distinct colors for adjacent vertices?). These problems give rise to many scheduling and routing problems with industrial importance...



 Now we define the problems that are addressed in this thesis and state the specific objectives of this research work.\\
 ....\\
 ....\\
 ....\\
 The objectives of this research work are:
 \begin{itemize}
    \item To ...  
    \item To ...
\end{itemize} 

\subsection{Major Contributions/ Findings}
The major contributions of this thesis are:
\begin{itemize}

    \item ...  
    \item ...
\end{itemize} 

\subsection{Organization of the Thesis}
Chapter-wise division of the results are presented here.

 \begin{enumerate}[label=\textbf{Chapter \arabic*},leftmargin=*]\setcounter{enumi}{0}

    \item provides a brief introduction to the two graph modification problems, namely, Subgraph Complementation and Seidel's Switching.  A survey on the related literature is also presented in this chapter.

\vspace{-0.2cm}
    \item   outlines the basic notations that are used throughout the thesis and defines the terms used in graph theory. In addition to this, we formally state the research problem and some folklore observations that will be useful in obtaining the results.
\vspace{-0.2cm}
    \item ....
\vspace{-0.2cm}
    \item ....
\vspace{-0.2cm}
    \item ....

\end{enumerate}

\section{Overview of the Thesis}
The scholar shall describe the relevant Methodology, Salient features, and Results of her/his research work.



\section{Conclusion}

Concluding remarks of the thesis are given here.
\newpage

\addcontentsline{toc}{section}{Details of Visible output from the Ph.D. work} 
\section*{Details of Visible output from the Ph.D. work}
\subsection*{Refereed Journals}
    \begin{itemize}
        \item D. G. Corneil, H. Lerchs, and L. S. Burlingham, “Complement reducible graphs,” Discret. Appl. Math., vol. 3, no. 3, pp. 163–174, 1981, \href{https://doi.org/10.1016/0166-218X(81)90013-5}{https://doi.org/10.1016/0166-218X(81)90013-5}  
    \end{itemize}

\subsection*{Refereed Conferences}

    \begin{itemize}     
        \item  Gartland and D. Lokshtanov, “Independent set on Pk-free graphs in quasi-polynomial time,” in 61st IEEE Annual Symposium on Foundations of Computer Science, FOCS 2020, Durham, NC, USA, November 16-19, 2020, S. Irani, Ed. IEEE, 2020, pp. 613–624, \href{https://doi.org/10.1109/FOCS46700.2020.00063}{https://doi.org/10.1109/FOCS46700.2020.00063}.
    \end{itemize}
\subsection*{Book Chapters}
\begin{itemize}     
        \item
    \end{itemize}
\subsection*{Patents granted}
     \begin{itemize}     
        \item
    \end{itemize}
\subsection*{Software copyrights granted}
     \begin{itemize}     
        \item
    \end{itemize}
\subsection*{Technology Transfer
executed}
     \begin{itemize}     
        \item
    \end{itemize}

 \newpage
\addcontentsline{toc}{section}{Selected References}
% \nocite{*}
\renewcommand\refname{Selected References}

\bibliographystyle{IEEEtran}
\bibliography{references.bib}


%\subsection*{Related Papers Published}


\end{document}
